\section{Fast Fourier Transform}
\label{Chap:FFT}
\begin{frame}[t]{Chapter \autoref{Chap:FFT}}
\tableofcontents[sections= \value{section}]
\end{frame}

\begin{frame}[t,allowframebreaks]{Introduction to FFT}

\begin{itemize}

\item The FFT is an \textcolor{red}{algorithm} which implements the DFT in an efficient manner


\item Put figure of comparison

	\begin{itemize}
	\item When $\text{log}_{10}(N) \, = \, 2$ which means that $N \, = \, 100$ samples, there is no a big difference between the DFT and the FFT
	
	\item As $\text{log}_{10}(N)$ increase, like at 6 (so $N$ is 1 million sample), the gap in the number of operation increase allot 
	\end{itemize}

\end{itemize}


\end{frame}

\begin{frame}[t,allowframebreaks]{FFT Algorithm details}

\begin{itemize}

\item We can compute 2 points of the DFT using only 1 multiplication

\item Many algorithms used in FFT, we look at the one so called \textcolor{red}{decimation in time}


\end{itemize}


\end{frame}
